% Part: applied-modal-logics
% Chapter: temporal-logic
% Section: possible-histories

\documentclass[../../../include/open-logic-section]{subfiles}

\begin{document}

\olfileid{aml}{tl}{poss}

\olsection{సాధ్య చరిత్రలు}

ఇప్పటివరకు వాడిన కాలిక తర్క సంబంధ నమూనాలు చాలా
సౌలభ్యమైనవి; ఎందుకంటే ప్రాప్యత సంబంధంపై ఎలాంటి
పరిమితులూ విధించాల్సిన అవసరం లేదు. అందువల్ల వాటిలో
గతం, భవిష్యత్తు రెండూ శాఖలుగా విడిపోవచ్చు. కానీ
ఘటనల అనుక్రమాలను సాధ్య చరిత్రలుగా చూసే, శాఖావిభజనకు
మరింత ``మోడల్'' భావనను కూడా పరిశీలించవచ్చు. దీనికి
మన భాషను తప్పనిసరిగా మార్చాల్సిన అవసరం లేదు. అయినా
``సాధారణ'' మోడల్ కారకాలు $\Box$,~$\Diamond$లను
చేర్చవచ్చు; కాలంతో పాటు కర్తల జ్ఞానంలో వచ్చే మార్పులను
సూచించడానికి జ్ఞానసంబంధ ప్రాప్యత సంబంధాలనూ చేర్చవచ్చు.

\begin{defn}
  కాలిక భాషకు \emph{సాధ్య చరిత్రల నమూనా} ఒక త్రయం
  $\mModel{M} = \tuple{T, C, V}$; ఇందులో
  \begin{enumerate}
    \item $T$ శూన్యం కాని సమితి; దీని సభ్యులను కాలంలోని
      స్థితులుగా అర్థం చేసుకుంటాం.
    \item $C$ ఒక వ్యవస్థ యొక్క గణన మార్గాల లేదా
      \emph{సాధ్య చరిత్రల} సమితి. అంటే, ప్రతి
      $s_i \in T$ అయ్యే స్థితులు $s_1$, $s_2$,~$s_3$,
      \dots{} గల అనుక్రమాల~$\sigma$ సమితి~$C$.
    \item $V$ అనేది ప్రతి \tetoken{ప్రతిజ్ఞావాక్య
      చరానికి}{!!{propositional
      variable}}~$p$ కాల బిందువుల సమితి~$V(p)$ను
      కేటాయించే ప్రమేయకం.
  \end{enumerate}
  సరళత కోసం, ఒక చరిత్ర~$C$లో ఉంటే దాని ప్రతి
  శేష అనుక్రమమూ దానిలోనే ఉంటుందని సాధారణంగా
  అనుకుందాం. ఉదాహరణకు, $s_1$, $s_2$,~$s_3$
  అనే అనుక్రమం~$C$లో ఉంటే, $s_2$, $s_3$ అనుక్రమం,
  అలాగే~$s_3$ కూడా అందులో ఉంటాయి. అనుక్రమం~$\sigma$లో
  $s_i$,~$s_j$ రెండు స్థితులు కనిపిస్తే, $i < j$
  అయినప్పుడు $s_i \prec_\sigma s_j$ అంటాం.
  $t \in V(p)$ అయితే $p$ అనేది~$t$ వద్ద
  \emph{సత్యం} అంటాం.
\end{defn}

సంబంధిత మార్పు ఒక్కటే: నమూనా~$\mModel M$లో కాల
బిందువు~$t$ వద్ద \tetoken{ఒక సూత్రం}{!!a{formula}}
సత్యమో కాదో నిర్ణయించేటప్పుడు, $t$ స్థితిగా కనిపించే
చరిత్ర~$\sigma$కు సంబంధించి నిర్ణయిస్తాం. అయితే
\tetoken{ప్రతిజ్ఞావాక్య చరాల}{!!{propositional
variable}s} లేదా సత్యప్రమేయ సంధాయకాల అర్థవిజ్ఞానాన్ని
మార్చాల్సిన అవసరం లేదు. వాటిలో దేనికీ~$\sigma$తో
సంబంధం లేదు కనుక అవి \olref[sem]{defn:tmodels}లో
ఉన్నట్లుగానే ఉంటాయి. కానీ ఈ చరిత్రలకు సంబంధించి
భవిష్యత్ కారకం~$\Ftemp$ను తిరిగి నిర్వచించి,
$\Diamond$ కారకాన్ని చేర్చుతాం.

\begin{defn}\ollabel{defn:phmodels}
  నమూనా~$\mModel M$లో \emph{$t, \sigma$ వద్ద
  \tetoken{ఒక సూత్రం}{!!a{formula}}~$!A$ సత్యం},
  సంకేతంగా $\mSat{M}{!A}[t, \sigma]$:
  \begin{enumerate}
  \item\ollabel{defn:sub:phmodels-f} \indcase{!A}{\Ftemp
    !B}{$\mSat{M}{\indfrm}[t, \sigma]$ అప్పుడే, అప్పుడు
    మాత్రమే $t \prec_\sigma t'$ అయ్యే ఏదో $t' \in T$కు
    $\mSat{M}{!B}[t', \sigma]$.}
  \item\ollabel{defn:sub:phmodels-diamond} \indcase{!A}{\Diamond
    !B}{$\mSat{M}{\indfrm}[t, \sigma]$ అప్పుడే, అప్పుడు
    మాత్రమే $t$ కనిపించే ఏదో $\sigma' \in C$కు
    $\mSat{M}{!B}[t, \sigma']$.}
  \end{enumerate}
\end{defn}

ఇతర కాలిక, మోడల్ కారకాలను కూడా ఇదే విధంగా
నిర్వచించవచ్చు. ఇప్పుడు కాలరూపం, మోడల్ సాధ్యత
కలిసిన వాదనలను వ్యక్తం చేయగలం. ఉదాహరణకు,
``$p$ జరగదు, కానీ జరిగి ఉండవచ్చు'' అనే వాదనను
$\lnot \Ftemp p \land \Diamond \Ftemp p$
సూత్రంతో సూచించవచ్చు. ఒక బిందువు, చరిత్ర వద్ద
$p$ తరువాతి స్థితిలో సత్యం కాకపోయినా, ప్రత్యామ్నాయ
చరిత్రలో $p$ భవిష్యత్తులో సత్యమయ్యే అవకాశం ఉంటే
ఈ సూత్రం అక్కడ నిలుస్తుంది.

\end{document}
